\chapter*{Introduction}
\markboth{Introduction}{Introduction}
\phantomsection
\addcontentsline{toc}{chapter}{Introduction}
\label{chap:introduction}

% Source PDF pp. 4--12 (Introduction, pp. 1--9).
It is now nearly 140 years since Riemann introduced \cite{Rie-EDL} the
concept of a ``local system'' on
\(\mathbb{P}^1-\{\text{a finite set of points}\}\). His idea was that one
could (and should) study the solutions of an \(n\)'th order linear
differential equation by studying the rank \(n\) local system (of its local
holomorphic solutions) to which it gave rise. Riemann knew that a rank
\(n\) local system on \(\mathbb{P}^1-\{m\text{ points}\}\) was ``nothing
more'' than a collection of \(m\) invertible matrices \(A_i\) in
\(\operatorname{GL}(n,\mathbb{C})\) which satisfy the matrix equation
\(A_1A_2\cdots A_m=(\mathrm{id}_n)\), such collections taken up to
simultaneous conjugation by a single element of
\(\operatorname{GL}(n,\mathbb{C})\). He also knew each individual \(A_i\)
was, up to \(\operatorname{GL}(n,\mathbb{C})\) conjugacy, just the effect
of analytic continuation along a small loop encircling the \(i\)'th missing
point.

His first application of these then revolutionary ideas was to study the
classical Gauss hypergeometric function \cite{Rie-SG}, which he did by
studying rank two local systems on \(\mathbb{P}^1-\{\text{three points}\}\).
His investigation was a stunning success, in large part because any such
(irreducible) local system is \emph{rigid} in the sense that it is determined
up to isomorphism as soon as one knows separately the individual conjugacy
classes of all its local monodromies. By exploiting this rigidity, Riemann
was able to recover Kummer's transformation theory of hypergeometric
functions ``almost without calculation'' \cite{Rie-APM}.

It soon became clear that Riemann had been ``lucky'', in the sense that the
most local systems are not rigid. For instance, rank two irreducible local
systems on \(\mathbb{P}^1-\{m\text{ points}\}\), all of whose local
monodromies are non-scalar, are rigid precisely for \(m=3\). And rank \(n\)
irreducible local systems on \(\mathbb{P}^1-\{\text{three points}\}\), each
of whose local monodromies has \(n\) distinct eigenvalues, are rigid
precisely for \(n=1\) and \(n=2\).

On the other hand, some of the best known classical functions are solutions
of differential equations whose local systems are rigid, including both of
the standard generalizations of the hypergeometric function, namely
\({}_nF_{n-1}\), which gives a rank \(n\) local system on
\(\mathbb{P}^1-\{0,1,\infty\}\), and the Pochhammer hypergeometric
functions, which give rank \(n\) local systems on
\(\mathbb{P}^1-\{n+1\text{ points}\}\).

In the classical literature, rigidity or its lack is expressed in terms of
the vanishing or nonvanishing of the ``number of accessory parameters''.
But the object whose rigidity is classically in question is not a rank
\(n\) local system on \(\mathbb{P}^1-\{m\text{ points}\}\), but rather an
\(n\)'th order linear differential equation with rational function
coefficients which has regular singularities at the \(m\) missing points,
and no other singularities. In practice, one assumes also that each of the
local monodromies has \(n\) distinct eigenvalues, expressed classically by
saying that no two exponents differ by integers. One then looks for the
most general \(n\)'th order linear differential equation with rational
function coefficients which has regular singularities at the \(m\) missing
points, no other singularities, and whose indicial polynomial at each
missing point is the same as for the equation we started with. (This game
with indicial polynomials makes sense for any equation with regular
singularities, but it is only a meaningful game in the case where each of
its local monodromies has \(n\) distinct eigenvalues, for only then can we
be sure that any equation with the same indicial polynomials automatically
has isomorphic local monodromies.) The ``number of parameters'' upon which
such an equation depends is called the ``number of accessory parameters'',
or the ``number of constants in excess'' \cite[\S20.4]{Ince}. For example,
in the classical literature one finds that for second order equations with
\(m\) regular singularities, the number of accessory parameters is
\(m-3\) \cite[top of page 506]{Ince}.

From a modern point of view, what corresponds precisely to a rank \(n\)
local system on \(\mathbb{P}^1-\{m\text{ points}\}\) is an \(n\times n\)
first order system of differential equations with regular singularities
at the named points (i.e., an algebraic vector bundle with integrable
connection on \(\mathbb{P}^1-\{m\text{ points}\}\) with regular
singularities at the \(m\) missing points) \cite{De-ED}. So what corresponds
to rigidity for a local system is the absence of deformations of the
corresponding \(n\times n\) system which preserve local monodromy. But in
the classical literature, the question of deforming such a system while
preserving its local monodromy does not seem to be addressed. Even, or
perhaps especially, if we start with an \(n\)'th order equation with
regular singularities, there is a priori a great difference between
deforming it as an equation and deforming it as a system (in both cases
preserving local monodromies).

For an irreducible local system \(\mathcal{F}\) on
\(\mathbb{P}^1-\{m\text{ points}\}\), there is a simple cohomological
invariant, \(\operatorname{rig}(\mathcal{F})\), the ``index of rigidity'',
which measures the rigidity or lack thereof of \(\mathcal{F}\), cf.\
\hyperref[chap:1]{Chapter~1}. One denotes by
\[
  j:\mathbb{P}^1-\{m\text{ points}\}\longrightarrow\mathbb{P}^1
\]
the inclusion, one forms the sheaf \(j_*\operatorname{End}(\mathcal{F})\) on
\(\mathbb{P}^1\), and one then defines \(\operatorname{rig}(\mathcal{F})\)
to be the Euler characteristic
\(\chi(\mathbb{P}^1,j_*\operatorname{End}(\mathcal{F}))\). Since
\(\mathcal{F}\) is irreducible, we have
\[
  \operatorname{rig}(\mathcal{F})
  :=2-h^1(\mathbb{P}^1,j_*\operatorname{End}(\mathcal{F})).
\]
One proves (\hyperref[thm:1.1.2]{Theorem~1.1.2}) that the irreducible local system
\(\mathcal{F}\) is rigid if and only if
\(\operatorname{rig}(\mathcal{F})=2\), i.e., if and only if
\(h^1(\mathbb{P}^1,j_*\operatorname{End}(\mathcal{F}))\) vanishes. So the
integer \(h^1(\mathbb{P}^1,j_*\operatorname{End}(\mathcal{F}))\) appears as a
cohomological analogue of the number of accessory parameters, at least in
the cases where that number was defined classically. In terms of the
corresponding \(n\times n\) system, this \(h^1\) is the ``number'' of its
deformations to systems having the same local monodromies. So in case we
start with an \(n\)'th order equation, we should expect this \(h^1\) to be
larger than the number of accessory parameters, since the \(h^1\) allows
deformations as system, while in computing the number of accessory
parameters we allow only deformations as equation.

In the case of a rank \(n\) irreducible local system \(\mathcal{F}\) which
arises from an \(n\)'th order equation, and each of whose local monodromies
has all distinct eigenvalues, one can separately compute both the number,
say \(\rho\), of accessory parameters and the number
\(h^1(\mathbb{P}^1,j_*\operatorname{End}(\mathcal{F}))\). One finds a doubling:
\[
  h^1(\mathbb{P}^1,j_*\operatorname{End}(\mathcal{F}))=2\rho.
\]
(Both sides come out to be \(2-[(2-m)n^2+mn]\), cf.\
\cite[pp.\ 127--128]{Forsythe} for the calculation of \(\rho\), and
\hyperref[chap:1]{Chapter~1} for the calculation of \(h^1\).) That the
\(h^1\) turns out to be even is not a surprise, because
\(H^1(\mathbb{P}^1,j_*\operatorname{End}(\mathcal{F}))\) carries a symplectic
autoduality. But why the underlying \(n\)'th order equation should have
``twice as many'' system-deformations as equation-deformations seems
entirely mysterious. It is as though the group
\(H^1(\mathbb{P}^1,j_*\operatorname{End}(\mathcal{F}))\) carried a weight one
Hodge structure, in such a way that the ``holomorphic part''
\(H^{1,0}\) corresponded to deformations as equation. But there is almost
nothing to back up such speculation.

Let us return to the consideration of a rank \(n\) local system on
\(\mathbb{P}^1-\{m\text{ points}\}\), in its incarnation as a collection,
taken up to simultaneous conjugation, of \(m\) elements \(A_i\) of
\(\operatorname{GL}(n,\mathbb{C})\) whose product is \(1\). Such a local
system is said to be irreducible if the subgroup of
\(\operatorname{GL}(n,\mathbb{C})\) generated by all the \(A_i\) is
irreducible, i.e., if there exists no proper nonzero subspace of
\(\mathbb{C}^n\) which is respected by all the \(A_i\). Given a local
system, we extract, at each of the \(m\) missing points, the conjugacy class
of its local monodromy at that point (concretely, the Jordan normal form
of each separate \(A_i\)), and call this the numerical data of our local
system. There is also the notion of abstract numerical data of rank \(n\)
on \(\mathbb{P}^1-\{m\text{ points}\}\): one specifies at each of the
\(m\) missing points a conjugacy class in
\(\operatorname{GL}(n,\mathbb{C})\), i.e., a Jordan normal form.

Two basic problems in the subject are:

\noindent\textbf{Irreducible Recognition Problem.}\quad
Given abstract numerical data of rank \(n\) on
\(\mathbb{P}^1-\{m\text{ points}\}\), determine if it is the numerical
data attached to an irreducible rank \(n\) local system on
\(\mathbb{P}^1-\{m\text{ points}\}\).

\noindent\textbf{Irreducible Construction Problem.}\quad
Given abstract numerical data of rank \(n\) on
\(\mathbb{P}^1-\{m\text{ points}\}\) which one is told arises from an
irreducible local system on \(\mathbb{P}^1-\{m\text{ points}\}\), construct
explicitly at least one such local system. Or construct all such local
systems.

Because the index of rigidity \(\operatorname{rig}(\mathcal{F})\) can be
expressed in terms of the underlying numerical data of \(\mathcal{F}\) by
universal formulas, one can pose these two problems separately for each of
the a priori possible values \(2,0,-2,-4,\ldots\) of the index of rigidity.
This book is devoted to the solution of these two problems in the special
case of rigid local systems, \(\operatorname{rig}(\mathcal{F})=2\).

The case of more general local systems remains entirely open. Already the
next simplest case, \(\operatorname{rig}(\mathcal{F})=0\), seems out of
reach. Also the analogues of these problems for reducible local systems
remain entirely open.

Another problem which remains open is this. Suppose we are given a rank
\(n\) irreducible rigid local system \(\mathcal{F}\) on
\(\mathbb{P}^1-\{m\text{ points}\}\). We know \cite{De-ED} that it is the
local system attached to a (unique, by rigidity) \(n\times n\) system on
\(\mathbb{P}^1-\{m\text{ points}\}\), with regular singular points at the
\(m\) missing points. Is this this \(n\times n\) system in fact (the system
attached to) an \(n\)'th order equation with regular singular points only
at the \(m\) missing points? By \cite{Ka-CV} any system has, Zariski
locally, cyclic vectors. So at the expense of allowing finitely many
additional ``apparent singularities'' in our \(n\)'th order equation, we
can always get one whose local system is \(\mathcal{F}\). The question is
whether there exists such an equation \emph{without} apparent
singularities. It should be remarked that if we drop the word ``rigid''
from this question, a ``counting constants'' argument of Poincar\'e
\cite[pp.\ 314--315 of Tome II, where he counts projective representations]{Poin}
suggests that ``in general'' it should \emph{not} be possible to avoid
apparent singularities in this ``strong form'' of the Riemann--Hilbert
problem. Another heuristic argument that one cannot avoid apparent
singularities, pointed out to me by Washnitzer, is this. Consider an
irreducible \(n\)'th order equation with regular singular points on
\(\mathbb{P}^1-\{m\text{ points}\}\), each of whose local monodromies has
\(n\) distinct eigenvalues. Suppose the underlying local system, say
\(\mathcal{F}\), is not rigid. Then the calculation \(h^1=2\rho\) discussed
above suggests that \(\mathcal{F}\) has a \(2\rho\)-dimensional deformation
space of local systems with the same local monodromies, and that only for
deformations \(\mathcal{G}\) in a \(\rho\)-dimensional subspace will there
be an \(n\)'th order equation on \(\mathbb{P}^1-\{m\text{ points}\}\),
i.e., without apparent singularities, with \(\mathcal{G}\) as monodromy.

Let us now turn to a more detailed discussion of the contents of this book.
Although the Irreducible Recognition Problem is, on its face, an elementary
problem about multiplying complex matrices which could be explained to a
bright high school student, our solution for rigids is, unfortunately,
far from elementary.

We begin with the trivial observation that on
\(\mathbb{P}^1-\{m\text{ points}\}\), any rank one local system is both
irreducible and rigid. Our basic idea is to construct two sorts
``operations'' (which we call ``middle convolution'' and ``middle tensor
product'') on a suitable collection of irreducible local systems which
preserve the index of rigidity and whose effect on local monodromy we can
calculate. The ``middle tensor'' operation offers no difficulty; all the
work is in working out the theory of middle convolution.

What does convolution have to do with rigid local systems? The idea is
simple. The earliest known, and still the best known, rigid local system
is the local system of solutions of the second order differential equation
% The supplied erratum replaces the erroneous x in the original coefficient.
\[
  \lambda(1-\lambda)\left(\frac{d}{d\lambda}\right)^2f
  +\bigl(c-(a+b+1)\lambda\bigr)\left(\frac{d}{d\lambda}\right)f-abf=0
\]
satisfied by the Gauss hypergeometric function \(F(a,b,c,\lambda)\).
A solution is given \cite[page 293]{WWW} by the integral
\[
  \int x^{a-c}(1-x)^{c-b-1}(\lambda-x)^{-a}\,dx.
\]
Our key observation is that formally, this integral is the additive
convolution
\[
  \int f(x)g(\lambda-x)\,dx
\]
of \(f(x):=x^{a-c}(1-x)^{c-b-1}\) and \(g(x):=x^{-a}\). We then view
\(f(x)\) as incarnating a rigid local system, and \(g(x)\) as incarnating
a Kummer sheaf on \(\mathbb{G}_m\), and think about forming the additive
convolution of two such objects. In some sense, our entire book consists
of first making sense of this, and then exploiting it.

In \hyperref[chap:2]{Chapter~2} we define the middle convolution operators,
and work out their basic properties. The theory of perverse sheaves is the
indispensable setting for this theory. The theory we need is that of middle
convolution on \(\mathbb{A}^1\) as additive group. However, we also devote
some attention to the theory on \(\mathbb{G}_m\), where it ties in nicely
with our previous work \cite{Ka-GKM} and \cite{Ka-ESDE} on Kloosterman and
hypergeometric sheaves and differential equations.

\hyperref[chap:3]{Chapters~3} and \hyperref[chap:4]{4} are devoted to the
proof that our middle convolution operators do in fact preserve the index
of rigidity, and to calculating their effect on local monodromy. Here the
main technical tool is the \(\ell\)-adic Fourier Transform in
characteristic \(p>0\). In \hyperref[chap:3]{Chapter~3}, we show that
Fourier Transform preserves the index of rigidity in characteristic \(p\).
Because middle convolution in characteristic \(p>0\) has a simple
expression in terms of Fourier Transform, we find that middle convolution
in characteristic \(p\) also preserves the index of rigidity. Because
Laumon has worked out the precise effect of Fourier Transform on local
monodromy, we also get the effect of middle convolution on local monodromy,
still in characteristic \(p>0\). In \hyperref[chap:4]{Chapter~4}, we use a
specialization argument to show that these results on middle convolution
still hold in characteristic zero (despite the fact that the Fourier
Transform no longer exists).

The next step is to show, in \hyperref[chap:5]{Chapter~5}, that any rigid
irreducible local system can be built up from a rank one local system by
applying a finite sequence of middle convolution and middle tensor
operations. The proof of this last step gives us an algorithm to calculate,
for any given irreducible rigid local system \(\mathcal{F}\), exactly what
sequence of operations to apply to what rank one local system
\(\mathcal{L}\) in order to end up with \(\mathcal{F}\). This algorithm is
thus a solution to the Irreducible Construction Problem for rigids. This
algorithm depends only on the ``numerical data'' of \(\mathcal{F}\).

Roughly speaking, we solve the Irreducible Recognition Problem (in
\hyperref[chap:6]{Chapter~6}) by showing that if we are given some abstract
numerical data which is rigid, then it comes from an irreducible
\(\mathcal{F}\) if and only if the algorithm, applied ``formally'', gives
a meaningful answer.

In \hyperref[chap:7]{Chapter~7}, we explore some of the diophantine aspects
of rigidity. In \hyperref[chap:8]{Chapter~8}, we reinterpret our
construction of rigids in terms of ``pieces'' of the relative de Rham
cohomology of suitable families of varieties. In
\hyperref[chap:9]{Chapter~9}, we use this cohomological expression,
together with an easy but previously overlooked generalization of our
earlier work on Grothendieck's \(p\)-curvature conjecture, to prove
Grothendieck's \(p\)-curvature conjecture for all those differential
equations on \(\mathbb{P}^1-\{m\text{ points}\}\) with regular singular
points whose underlying local systems are rigid and irreducible.

Let us now discuss what is \emph{not} done in this book. One could try to
classify systematically all rigid irreducible local systems, since one has
an algorithm to recognize their numerical data. Even a cursory glance at
the kinds of local monodromy one can get by starting with a rank one local
system and applying a cleverly chosen sequence of middle convolution and
middle tensor operations leaves one with the impression that there is a
fascinating bestiary waiting to be compiled.

One could also study the identities between special functions which
presumably result whenever a rigid irreducible \(\mathcal{F}\) can be built
in two or more different ways out of rank one local systems by successive
middle convolution and middle tensor operations. Already for the case of
\({}_nF_{n-1}\), any \(n\geq2\), there are in general a plethora of such
building paths. Do we get anything about \({}_nF_{n-1}\) which is not
already in the classical literature?

In characteristic zero, we have shown how to construct all rigid
irreducible local systems out of rank one local systems, by using the
operations of middle convolution and middle tensor. Our arguments in fact
begin in characteristic \(p>0\), where we prove the same result for rigid
irreducible local systems which are everywhere tamely ramified. But in
characteristic \(p\), the everywhere tame local systems are by far the
least interesting ones. What can be said about arbitrary rigid irreducible
local systems in characteristic \(p\)? Is it true that any rigid
irreducible local system in characteristic \(p\) on
\(\mathbb{P}^1-\{m\text{ points}\}\) is built out of a rank one local
system by finitely iterating the operations Fourier Transform, middle
tensor with a rank one local system, and pullback by an automorphism of
\(\mathbb{P}^1\)? For example, all the irreducible \(\ell\)-adic
hypergeometrics are rigid, and they are all obtained in this way
\cite[proof of 8.5.3]{Ka-ESDE}.

In characteristic zero, the analogue of a not necessarily tame local
system is a differential equation which does not necessarily have regular
singular points. There is no difficulty in defining the index of rigidity
in the holonomic \(\mathcal{D}\)-module context, cf.\
\cite[3.7.3 and 2.9.8.1]{Ka-ESDE}. One knows, for example, that the
generalized hypergeometric equations studied in \cite{Ka-ESDE} are rigid.
One also has the \(\mathcal{D}\)-module Fourier Transform. It \emph{should}
be true that Fourier Transform preserves the index of rigidity in the
\(\mathcal{D}\)-module context, but this is unknown. The main stumbling
block to proving this is the absence of an \(\mathcal{D}\)-module analogue
of Laumon's theory of local Fourier Transform and his stationary phase
theorem relating the local and global Fourier Transforms: Laumon's theory
in the \(\ell\)-adic case was the main technical tool in our proof that
Fourier Transform preserves index of rigidity. If the
\(\mathcal{D}\)-module analogue of Laumon's local Fourier Transform exists
and satisfies stationary phase, one can then use stationary phase to
``compute'' what the local Fourier Transforms must be, in terms of slope
decompositions at \(\infty\) of global Fourier Transforms of suitable
``canonical extensions'' in the sense of \cite[\S2.4.11]{Ka-DGG} of various
completions of the input \(\mathcal{D}\)-module. In this sense, one could
say that the theory of local Fourier Transform does already exist, and
that ``all'' that one lacks is the \(\mathcal{D}\)-module analogue of
stationary phase.

If one could prove that the \(\mathcal{D}\)-module Fourier Transform
preserves index of rigidity, or even that it preserves rigid objects, one
could ask if any irreducible \(\mathcal{D}\)-module is built out of an
irreducible \(\mathcal{D}\)-module of generic rank one by finitely iterating
the operations Fourier Transform, middle tensor with an irreducible
object of generic rank one, and pullback by an automorphism of
\(\mathbb{P}^1\).

Another question we are unable to treat is the following. Given a rigid
irreducible local system, say \(\mathcal{F}\), of rank \(n\), consider its
geometric monodromy group \(G_{\mathrm{geom}}\), defined as the algebraic
subgroup of \(\operatorname{GL}(n,\mathbb{C})\) which is the Zariski closure
of the monodromy representation which \(\mathcal{F}\) ``is''. This group
\(G_{\mathrm{geom}}\), and consequently its Lie algebra, is determined up
to \(\operatorname{GL}(n,\mathbb{C})\)-conjugacy by \(\mathcal{F}\), and
hence by the numerical data of \(\mathcal{F}\). How can we determine
\(G_{\mathrm{geom}}\), or even its Lie algebra, as a function of the
numerical data. Even in asking when \(G_{\mathrm{geom}}\) is finite, which
we prove is equivalent to (the associated differential equation's) having
\(p\)-curvature zero for almost all primes \(p\), we do not know how to
read this from the numerical data.

There is also a nagging technical point that is not treated in this book.
In \hyperref[chap:1]{Chapter~1}, we show (\hyperref[thm:1.1.2]{Theorem~1.1.2}) that
for complex irreducible local systems on
\(\mathbb{P}^1-\{m\text{ points}\}\) over \(\mathbb{C}\), rigidity is
equivalent to having index of rigidity equal to \(2\). And we show
(\hyperref[thm:5.0.2]{Theorem~5.0.2}) that in any characteristic \(\ne\ell\),
having index of rigidity equal to \(2\) is a sufficient condition for an
irreducible \(\ell\)-adic local system to be rigid. However, we do not
show, even over \(\mathbb{C}\), that for irreducible \(\ell\)-adic local
systems, being rigid is in fact equivalent to having index of rigidity
equal to \(2\).

It is perhaps striking that although this book is concerned with problems
that go back to Riemann, it depends for its very existence on a great deal
of mathematics that did not exist until quite recently: Grothendieck's
\'etale cohomology theory \cite{SGA}, Deligne's proof of his far-reaching
generalization of the original Weil Conjectures \cite{De-WeilII}, the
theory of perverse sheaves \cite{BBD}, Laumon's work on the \(\ell\)-adic
Fourier Transform \cite{Lau-TF}, all these are indispensable ingredients.

My interest in rigid local systems was first aroused by a conversation
with Ofer Gabber some years ago, who told me about some lectures Deligne
had given at I.H.E.S.\ on them. It was later re-aroused by conversations
with Carlos Simpson, who was pursuing them from quite a different point
of view than that used here. It is a pleasure to acknowledge helpful
discussions with Deligne, Gabber and Simpson, and to thank Beilinson and
Faltings for asking some incisive questions. I would also like to thank
the referee, whose helpful comments and suggestions led to a number of
corrections to and clarifications of the original manuscript.

I gave a series of lectures on some of the material in this book in
January, 1991 at the University of Minnesota as an Ordway Visitor. The
material on middle convolution was presented in March of 1993 at Johns
Hopkins, at a Symposium in honor of Professor Igusa. I also gave lectures
on some of this book in May of 1993 in Berkeley as a Miller Visiting
Fellow. It is a pleasure to thank all of those institutions for their
support and hospitality.
